\subsection{Montel coordinate representation}
\label{sec:montel-coord-represents}

This subtree isolates the \code{coord\_represents\_uniformization} field of
\code{GenusZeroNormalizedMontelPatchSelector}.  It is the local statement that
each selected Montel coordinate represents the global uniformization in the
assigned public chart of \(\code{OnePoint}\,\C\).

\begin{lemma}[Montel selector coordinate represents the uniformization (open)]
\label{lem:montel-selector-coord-represents-uniformization}
\uses{lem:montel-selector-source-patch-cover, lem:global-gluing-coord-mem-target, lem:montel-selector-chart-inverse-local-value, lem:montel-selector-normal-limit-agrees-with-uniformization, lem:montel-selector-overlap-compatible-normalized-values}
For every selected source patch \(U_i\), source point \(x\in U_i\), and assigned
target chart \(\psi_i\) on \(\code{OnePoint}\,\C\), the selected local coordinate
represents the global uniformization:
\[
  \psi_i^{-1}(z_i(x)) = F(x).
\]
This is the field
\code{GenusZeroNormalizedMontelPatchSelector.coord\_represents\_uniformization}.
The Lean projection exists as a structure field, but the construction of a
selector carrying this field is the open Montel assembly obligation, so this
blueprint node intentionally carries no \code{\textbackslash lean\{\}} citation.
\end{lemma}

\begin{lemma}[Selected source patches cover the surface]
\label{lem:montel-selector-source-patch-cover}
The finite global patch-family structure supplies a selected source patch
containing each point of \(X\).  Its \code{patch\_cover} field is the cover
witness used when turning local coordinate statements into a global formula.
\lean{JacobianChallenge.HolomorphicForms.GenusZeroGlobalPatchFamily}
\leanok
\end{lemma}
\begin{proof}\leanok\end{proof}

\begin{lemma}[Selected chart inverse gives the local candidate value (open)]
\label{lem:montel-selector-chart-inverse-local-value}
\uses{lem:global-gluing-coord-mem-target, lem:mathlib-open-partial-homeomorph-map-target, lem:mathlib-open-partial-homeomorph-right-inv}
On a selected source patch, the local coordinate lies in the assigned target
chart.  Applying the target-chart inverse therefore gives the local candidate
point of \(\code{OnePoint}\,\C\) that should become the global uniformization
value.

This is a planning node for the project-specific reconstruction package.  Its
formal floor is already visible: target membership comes from the global patch
family, and Mathlib's open-partial-homeomorphism API says that the chart inverse
lands in the chart source and reconstructs the original coordinate.
\end{lemma}

\begin{lemma}[Mathlib: target points map back to the source]
\label{lem:mathlib-open-partial-homeomorph-map-target}
For an open partial homeomorphism \(e\), if \(z\in e.\mathrm{target}\), then
\(e^{-1}(z)\in e.\mathrm{source}\).
\lean{OpenPartialHomeomorph.map_target}
\leanok
\end{lemma}
\begin{proof}\leanok\end{proof}

\begin{lemma}[Mathlib: chart inverse reconstructs target coordinates]
\label{lem:mathlib-open-partial-homeomorph-right-inv}
For an open partial homeomorphism \(e\), if \(z\in e.\mathrm{target}\), then
\[
  e(e^{-1}(z)) = z.
\]
\lean{OpenPartialHomeomorph.right_inv}
\leanok
\end{lemma}
\begin{proof}\leanok\end{proof}

\begin{lemma}[Normalized Montel limit agrees with the uniformization (open)]
\label{lem:montel-selector-normal-limit-agrees-with-uniformization}
\uses{lem:montel-selector-source-patch-cover, lem:montel-selector-chart-inverse-local-value}
The normal-family/Montel extraction must choose local coordinates whose
target-chart inverse agrees with the global homeomorphism \(F:X\simeq
\code{OnePoint}\,\C\) on the selected source patch.

This is the analytic heart of the coordinate-representation field: it is not a
formal consequence of the existing gluing API, and it should remain an open
provider until the Montel patch-selector construction supplies it.
\end{lemma}

\begin{lemma}[Normalized local values are overlap-compatible (open)]
\label{lem:montel-selector-overlap-compatible-normalized-values}
\uses{lem:global-gluing-coord-mem-target}
On overlaps of selected source patches, the two target-chart inverse formulas
for the normalized Montel coordinates determine the same point of
\(\code{OnePoint}\,\C\).  This compatibility is what lets the local
representations define one global uniformization rather than unrelated chart
values.

This is a planning node with no current Lean declaration.  Downstream green
gluing lemmas consume this compatibility once the selector field is available,
but they should not be used as fake providers for constructing the field.
\end{lemma}
